\subsection{Vanishing of the three-torus obstruction}

For the standard real PL torus $T^3=\mathbb R^3/\mathbb Z^3$,
the obstruction of
Theorem~\ref{thm:natural-three-concordance-obstruction}
has the following value.

\begin{theorem}\label{thm:three-torus-concordance-obstruction}
With the ordinary bottom-fixed, top-free concordance convention,
\begin{equation}\label{eq:three-torus-obstruction-zero}
 c_{T^3}:\pi_0P_{\mathrm{TOP}}(T^3)
                   \longrightarrow H^2(T^3;\mathbb Z/2)
 \qquad\text{is zero.}
\end{equation}
\end{theorem}

\begin{proof}
Write $X=T^3\times I$.  We first make both end maps the
identity, and then construct a relative five-dimensional
structure whose normal invariant detects the given obstruction.

\emph{End maps and collars.}
For any $H\in P_{\mathrm{TOP}}(T^3)$, projecting $H(x,s)$
to $T^3$ gives a homotopy from the identity to its top map.
The endpoint normalization in
Theorem~\ref{thm:natural-three-concordance-obstruction}
permits us to assume that this top map $a$ is PL, without
changing $c_{T^3}([H])$.  The three-torus is irreducible and
contains an incompressible coordinate two-torus.  Hence
Waldhausen's homotopy--isotopy theorem applies to $a$:
$a$ is TOP isotopic to the identity
\cite[Theorem~7.1, pp.~80--81; proof pp.~81--85]{Waldhausen1968}.
The component injectivity in
\eqref{eq:three-manifold-full-automorphism-comparison}
therefore supplies a joint PL isotopy $\gamma_s$ from $a$
to the identity.  Its trace
\[
 E_\gamma(x,s)=(\gamma_s a^{-1}(x),s)
\]
is a PL concordance with bottom map the identity and top
map $a^{-1}$.  Replacing $H$ by $E_\gamma H$ makes both
end maps the identity and leaves $c_{T^3}$ unchanged.

We can moreover arrange that the resulting map $F$ is the
identity on product collars of both ends.  Explicitly, for
$0<r<1/2$ put $e_r(x,s)=(x,r+(1-2r)s)$, use
$e_rFe_r^{-1}$ on $T^3\times[r,1-r]$, and use the identity
on the two complementary collars.  At $r=0$ use $F$.
The pieces agree on their interfaces; applying the same
construction to $F^{-1}$ gives the inverse.  Uniform
continuity of $F$ and $F^{-1}$ on the compact cylinder
proves continuity at $r=0$.  This is an ordinary isotopy
relative to both ends, so it preserves the obstruction.

\emph{Removing the lift displacement.}
Since the bottom inclusion induces an isomorphism on
fundamental groups and $F$ fixes that bottom, $F_*$ is
the identity on $\mathbb Z^3$.  Choose the lift
\[
 \widetilde F(x,s)=(u(x,s),v(x,s)):
                    \mathbb R^3\times I\longrightarrow
                    \mathbb R^3\times I
\]
whose bottom restriction is $(x,0)$.  It is
$\mathbb Z^3$-equivariant.  On the upper collar its first
coordinate is $x+z$ for a constant $z\in\mathbb Z^3$,
whereas on the lower collar it is $x$.
Choose a PL function $\rho:I\to I$ that is zero on a
lower collar and one on an upper collar, and set
\begin{equation}\label{eq:three-torus-lift-correction}
 L_z([x],s)=([x-\rho(s)z],s).
\end{equation}
This is a PL homeomorphism: on the finitely many intervals
where $\rho$ is affine it is induced by an affine map,
and a finite subdivision accommodates the quotient
$\mathbb R^3\to T^3$.  It is the identity on both end
collars.  Its obstruction is zero.  The chosen lift of
$L_zF$ has first coordinate $u(x,s)-\rho(v(x,s))z$,
so it agrees with $x$ on both collars.  Replace $F$ by
$L_zF$ and retain the notation $F$.  We now have
\[
 \widetilde F(x,s)=(x+w(x,s),v(x,s)),
\]
where $w$ is $\mathbb Z^3$-periodic and $w=0$, $v=s$
on both end collars.  Consequently
\begin{equation}\label{eq:three-torus-relative-homotopy}
 k_t([x],s)=
 \bigl([x+t w(x,s)],(1-t)s+t v(x,s)\bigr)
\end{equation}
is a well-defined homotopy of maps from the identity to
$F$, relative to those collars.  The intermediate maps
need not be homeomorphisms.  Reparametrize $t$ so that
$k_t$ is constant near $t=0$ and near $t=1$.

\emph{The specified relative structure.}
Let $M_F$ be the mapping cylinder of $F:X\to X$, with
upper label $y=F(x)$, and put $Z=X\times I$.
Define
\begin{equation}\label{eq:three-torus-relative-structure}
 U:M_F\longrightarrow Z,
 \qquad U([x,t])=(k_t(x),t).
\end{equation}
On the upper boundary, $U$ is the identity in the
$y$-label; it is also the identity on the bottom and
side boundary.  The constant end portions of $k_t$ and
the collars on which $F$ is the identity make these
identifications products on neighborhoods of all faces,
including their intersections.  Thus $U$ preserves the
specified PL boundary reductions.  Both bottom inclusions
$X\to M_F$ and $X\to Z$ are homotopy equivalences, and
$U$ restricts to the identity on that bottom.  Hence
$U$ is a homotopy equivalence of compact five-manifolds
relative to their full boundary.

Transport the stable TOP normal bundle of $M_F$ across
$U$, together with the induced identification of its
Spivak normal fibration and the fixed boundary comparison.
This gives the relative normal invariant
\[
 \nu(U)\in[Z,\partial Z;G/\mathrm{TOP},*].
\]
Its surgery obstruction is zero because $U$ is a
homotopy equivalence.  The equality
$\operatorname{Wh}(\mathbb Z^3)=0$
\cite[Example~24.16, p.~278]{Ranicki1992TopologicalManifolds}
also identifies the simple and homotopy decorations here.

\emph{The exact assembly input.}
The relative normal-invariant sequence applies to the
five-dimensional manifold $Z$.  In its normal-invariant
term, Poincar\'e--Lefschetz duality identifies the surgery
obstruction map with quadratic assembly:
\begin{equation}\label{eq:three-torus-relative-assembly}
 \begin{split}
 [Z,\partial Z;G/\mathrm{TOP},*]
   &\cong H_5(Z;\mathbb L\langle1\rangle)
    \cong H_5(T^3;\mathbb L\langle1\rangle),\\
 \theta&=A:
 H_5(T^3;\mathbb L\langle1\rangle)
       \longrightarrow L_5(\mathbb Z[\mathbb Z^3]).
 \end{split}
\end{equation}
Here $\mathbb L\langle q\rangle$ is the $q$-connective
quadratic surgery spectrum.  The precise relative input
is \cite[Definition~18.2, pp.~197--198; the relative-version
paragraph on p.~198; Proposition~18.3 and Theorem~18.5,
pp.~198--201]{Ranicki1992TopologicalManifolds}.
The free-abelian assembly calculation in
\cite[Example~24.16, p.~278]{Ranicki1992TopologicalManifolds}
gives
\[
 H_5(T^3;\mathbb L\langle0\rangle)
                \xrightarrow{\ \cong\ }
 L_5(\mathbb Z[\mathbb Z^3]),
\]
since five is at least the rank three.
The cofiber of
$\mathbb L\langle1\rangle\to\mathbb L\langle0\rangle$
is $H\mathbb Z$.  Its exact sequence contains
\[
 H_6(T^3;\mathbb Z)\longrightarrow
 H_5(T^3;\mathbb L\langle1\rangle)\longrightarrow
 H_5(T^3;\mathbb L\langle0\rangle)\longrightarrow
 H_5(T^3;\mathbb Z).
\]
The outside terms vanish, so the middle comparison is
an isomorphism.  Thus $\theta$ in
\eqref{eq:three-torus-relative-assembly} is injective,
and $\nu(U)=0$.

\emph{Returning to the relative PL obstruction.}
Forgetting the spherical-fibration comparison in this
zero normal invariant gives a homotopy, relative to the
specified boundary comparison, from the transported
TOP normal reduction of $M_F$ to the standard normal
reduction of $Z$.  The latter has its standard PL lift.
Pulling that lift back along the homotopy and then along
$U$ gives a stable PL normal reduction on $M_F$ extending
its prescribed boundary reduction.  In particular its
relative normal Kirby--Siebenmann obstruction is zero.
The stable tangent and normal obstructions agree, since
their Whitney sum is trivial and the obstruction is
additive with mod-two coefficients.  Therefore
\[
 \operatorname{ks}(M_F,\partial M_F)=0.
\]
By the defining mapping-cylinder suspension of the
relative Casson--Sullivan class, followed by the interval
suspension $H^2(T^3;\mathbb Z/2)\cong
H^3(X,\partial X;\mathbb Z/2)$, this implies
$c_{T^3}([F])=0$.  Each preceding normalization preserved
$c_{T^3}$, which proves
\eqref{eq:three-torus-obstruction-zero}.
\end{proof}

The conclusion determines this particular component
obstruction.  A classification of the ordinary concordance
components would require further information beyond its value.

\paragraph{A proof by finite covers.}
There is also a direct covering argument for
Theorem~\ref{thm:three-torus-concordance-obstruction}.  Its geometric
input is the relative local deformation theorem of Edwards--Kirby
\cite[Theorem~5.1(1)--(2) and its boundary-pair Note, p.~70;
proof of Corollary~1.1, p.~77]{EdwardsKirby1971}.
We give the argument because it also describes the ordinary components
of sufficiently large lifts.

Put $B=T^3$, $M=B\times I$, and
$\mathcal P=P_{\mathrm{TOP}}(B)$.  For $N\geq1$ let
$p_N([x])=[Nx]$ and $q_N=p_N\times1_I$.
Every $H\in\mathcal P$ has a unique lift $H_N$ through $q_N$
which fixes the bottom pointwise.  Indeed $H_*=1$ on
$\pi_1M$, because the bottom inclusion is a homotopy equivalence.
The bottom-fixed universal lift and its inverse are therefore
$\mathbb Z^3$-equivariant homeomorphisms.  Write the lift as
\[
 \widetilde H(x,t)=(x+u(x,t),v(x,t)).
\]
The function $u$ is periodic and hence bounded on the compact
fundamental cylinder.  If its sup norm is at most $C$, then
\begin{equation}\label{eq:three-torus-cover-displacement}
 H_N([x],t)=
 \left([x+N^{-1}u(Nx,t)],v(Nx,t)\right),
 \qquad
 \sup_{x,t}d_B(\operatorname{pr}_BH_N(x,t),x)\leq C/N,
\end{equation}
where $d_B$ is the quotient sup metric.
The same bound holds for the top map $g_N$.

For a bottom-fixed homeomorphism $J$ with top map $j$, and $0<h\leq1$,
define
\[
 T_h(J)(x,t)=
 \begin{cases}
  (\operatorname{pr}_BJ(x,t/h),
                h\operatorname{pr}_IJ(x,t/h)),&0\leq t\leq h,\\
  (j(x),t),&h\leq t\leq1.
 \end{cases}
\]
The formulas agree at $t=h$.  The lower piece is conjugation
into the lower slab; the upper piece is a product map.
Replacing $J,j$ by $J^{-1},j^{-1}$ gives the inverse.
Thus $T_h(J)\in\mathcal P$ and $T_1(J)=J$.
For every fixed $h_0>0$ the formulas and their inverses are jointly
continuous for $h\in[h_0,1]$, by pasting on the two closed regions.
They give an ordinary path from $J$ to $T_{h_0}(J)$.
With the product max metric on $M$,
\[
 \sup_{z\in M}d_M(T_h(H_N)(z),z)\leq\max\{C/N,h\}.
\]
In the lower slab both time coordinates lie in $[0,h]$; in the
upper slab time is unchanged, and the horizontal estimate holds
on both pieces.

Edwards--Kirby's theorem gives a neighborhood $\mathcal U$ of the
identity whose bottom-fixed members contract to the identity through
bottom-fixed homeomorphisms.  To see the exact relative hypothesis,
take $C=U=M$ in that theorem and take
$D_1=V_1=B\times\{0\}$ in its boundary-pair Note.  The bottom
is both closed and open in $\partial M$, so $V_1$ is a neighborhood
of $D_1$ there.  For compact $M$, the proper self-embedding space
in the theorem is $\operatorname{Homeo}(M)$, as in its cited
Corollary~1.1 proof.  The deformation ends at the identity on all
of $M$ and preserves the prescribed bottom.
Choose $\varepsilon>0$ so that uniform displacement less than
$\varepsilon$ implies membership in $\mathcal U$, and fix
$0<h_0<\min\{1,\varepsilon\}$.  For every sufficiently large $N$,
the displayed estimate puts $T_{h_0}(H_N)$ in $\mathcal U$.
The compression path followed by the local deformation proves
\begin{equation}\label{eq:three-torus-large-cover-component}
 [H_N]=[1]\quad\text{in }\pi_0\mathcal P
 \quad\text{for every sufficiently large }N.
\end{equation}

It remains to verify the covering identity for the obstruction,
including its endpoint data.  Let $g$ be the top map of $H$.
The three-dimensional automorphism comparison, followed by
Waldhausen's theorem as above, gives an ordinary isotopy
$\beta_s$ from $g$ to the identity.  The maps
\[
 A_s(x,t)=(\beta_{st}g^{-1}(x),t),\qquad H_s=A_sH
\]
form a bottom-fixed isotopy from $H$ to a map $F=H_1$ fixing
both ends.  Their inverse formulas use $g\beta_{st}^{-1}$,
so the total inverse family is continuous as well.
Lift this entire isotopy through $q_N$, beginning at $H_N$.
Its endpoint $F_N$ fixes the bottom but can have upper map
translation by $a/N$, for some $a\in\mathbb Z^3$.

For any fixed $w\in\mathbb R^3$ put
$S_w([x],t)=([x-tw],t)$.  This is PL on the whole cylinder:
its universal lift is affine, and subdivision along the finitely
many quotient hyperplanes in a fundamental cube gives a finite
PL description.  Its inverse is $S_{-w}$, and it fixes the bottom.
Set
\[
 K_N=S_{a/N}F_N,\qquad K=S_aF.
\]
Both maps fix both ends, because $a$ is integral, and
\[
 q_NS_{a/N}=S_aq_N,\qquad q_NK_N=Kq_N.
\]
The obstruction is a component homomorphism and vanishes on PL
concordances.  Therefore
\[
 c_B[K_N]=c_B[H_N],\qquad c_B[K]=c_B[H].
\]

For the both-ends-fixed maps $K_N,K$, their marked auxiliary
mapping cylinders form a covering.  The tangent microbundle,
its PL boundary reductions, and the comparisons induced by these
maps all pull back: $q_N$ is locally PL and the displayed square
commutes.  Naturality of the relative obstruction in the definition
of $c_B$ gives
\[
 \operatorname{cs}_{\mathrm{PL}}(K_N)
       =q_N^*\operatorname{cs}_{\mathrm{PL}}(K)
       \quad\text{in }H^3(M,\partial M;\mathbb Z/2).
\]
The interval suspension $\sigma:H^2(B;\mathbb Z/2)\to
H^3(M,\partial M;\mathbb Z/2)$ satisfies
$q_N^*\sigma=\sigma p_N^*$, since $q_N$ is the identity
on the interval.  The defining both-ends-fixed formula for $c_B$
now proves
\begin{equation}\label{eq:three-torus-cover-pullback}
 c_B[H_N]=p_N^*c_B[H].
\end{equation}
This establishes the covering formula with its required markings.

Finally, if $x_1,x_2,x_3$ are the circle generators in mod-two
degree-one cohomology, then $p_N^*(x_ix_j)=N^2x_ix_j$.
Thus every odd $N$ acts identically on $H^2(T^3;\mathbb Z/2)$.
Choose such an $N$ large enough for
\eqref{eq:three-torus-large-cover-component}; then
\eqref{eq:three-torus-cover-pullback} gives $c_B[H]=0$.
The contraction of $H_N$ need not be equivariant under the deck
group, so this argument makes no assertion about an isotopy of
the original concordance or surjectivity of PL comparison.
